\begin{helper}
Fix \(I\), a finite set \(R\) of terminal coordinates, and samples
\(\omega,\omega'\).  If \(R\) is exactly
\(\noderef{TwoBitePreTerminalRecordedPairs}(\omega,\varepsilon,I)\), if
\(\omega\) and \(\omega'\) have the same injection, and if they have the same
\(\noderef{TwoBiteEdgeCoordinateValue}\)-truth value on every coordinate of
\(R\), then \(R\) is also exactly
\(\noderef{TwoBitePreTerminalRecordedPairs}(\omega',\varepsilon,I)\).
\end{helper}
\begin{proof}
Write \(R\) for the recorded set in the statement.  By hypothesis,
\[
e\in R \quad\Longleftrightarrow\quad
e\in\noderef{TwoBitePreTerminalRecordedPairs}(\omega,\varepsilon,I)
\]
for every terminal coordinate \(e\).  We compare the defining predicates for
\(\omega\) and \(\omega'\).

First, projection containment is the same in the two samples.  It is defined
only from the injection and the red or blue projection of \(I\), and the
injections agree by hypothesis.  The fibers over base vertices are also the
same, so \(F_1\) agrees for every base vertex.  For \(F_2\), fix a base vertex
\(x\).  The finite list of possible neighbors \(y\) is the same, and the
\(F_1\)-status of each \(y\) is already the same.  If \(x\) and \(y\) have
opposite colors, \(\noderef{TwoBiteBaseAdj}\) is false in both samples.  Now
suppose \(x\) and \(y\) have the same color.  If their projected endpoints are
equal, then the relevant adjacency is a loop edge in the red graph or in the
blue graph, respectively.  Both graphs are simple graphs, so loop adjacency is
false in \(\omega\) and in \(\omega'\).  Hence the two
\(\noderef{TwoBiteBaseAdj}\)-truth values agree in this endpoint-equal case,
and no terminal coordinate is needed.

It remains to handle the same-color endpoint-distinct case.  Let \(q\) be the
increasing ordering of the two projected endpoints.  Whenever this adjacency
can affect the \(F_2\)-count, the endpoint \(y\) is in \(F_1\) and \(x\) is
projection-contained.  Therefore
\(\noderef{FixedSetPreTerminalF2QueriesRecorded}\), applied to the increasing
coordinate \(q\), puts the tagged coordinate for \(q\) in
\(\noderef{TwoBitePreTerminalRecordedPairs}(\omega,\varepsilon,I)\), hence in
\(R\).  The truth-table agreement on \(R\) then gives the same value of
\(\noderef{TwoBiteBaseAdj}(\omega,x,y)\) and
\(\noderef{TwoBiteBaseAdj}(\omega',x,y)\).  Thus the filtered set counted in
\(\noderef{TwoBiteStageF2}\) is the same for \(\omega\) and \(\omega'\), and
\(\noderef{TwoBiteStageF2}\) agrees.  Since
\(\noderef{TwoBiteStageRevealedVertex}\) is the disjunction of the
not-projection-contained, \(F_1\), and \(F_2\) alternatives, staged-revealed
vertices agree in the two samples.

We next prove that the \(X_x(I)\)-sets agree for every staged-revealed base
vertex \(x\).  Fix \(v\in I\).  Membership of \(v\) in
\(\noderef{TwoBiteX}(\omega,I,x)\) means that \(v\) is in the lifted
neighborhood of \(x\).  In the red case this is the red graph edge from the
red endpoint of \(x\) to \(\noderef{TwoBiteRedProjection}(\omega,v)\); in the
blue case it is the analogous blue edge.  The projected endpoint is the same
for \(\omega\) and \(\omega'\) because their injections agree.  If this
projected endpoint is equal to the endpoint of \(x\), then the lifted
neighborhood test is again a loop adjacency in the corresponding simple graph,
so it is false in both samples.  Thus membership agrees in this case without
using a recorded coordinate.

If the projected endpoint is distinct from the endpoint of \(x\), take the
increasing ordered coordinate between these two endpoints, with the tag
determined by their common color.  Since \(x\) is staged-revealed and the
projected endpoint of \(v\in I\) is projection-contained, the incident branch
of \(\noderef{TwoBitePreTerminalRecordedPairs}(\omega,\varepsilon,I)\) puts
this tagged coordinate in \(R\).  Agreement on \(R\) gives the same edge value
in \(\omega\) and \(\omega'\).  Hence membership in \(X_x(I)\) is the same in
the two samples for every \(v\in I\).  Therefore the projected pair sets
\((\noderef{TwoBiteX}(\omega,I,x)).\mathrm{image}\) and
\((\noderef{TwoBiteX}(\omega',I,x)).\mathrm{image}\) coincide in the
appropriate color.

Now unfold
\(\noderef{TwoBitePreTerminalRecordedPairs}(\omega',\varepsilon,I)\).  Its
outer terminal-coordinate condition is the same fixed terminal universe.  For
an incident-branch coordinate, the relevant staged-revealed and
projection-contained predicates for \(\omega'\) agree with those for
\(\omega\), so the coordinate satisfies the same incident branch for
\(\omega\), and hence lies in \(R\); conversely every incident-branch
coordinate for \(\omega\) satisfies the corresponding branch for \(\omega'\).
For an \(X_x(I)\)-pair coordinate, staged-revealed vertices agree and the
\(X_x(I)\)-sets just proved equal, so the \(X_x(I)\)-pair branch also holds
for \(\omega'\) exactly when it holds for \(\omega\).  Thus every terminal
coordinate belongs to
\(\noderef{TwoBitePreTerminalRecordedPairs}(\omega',\varepsilon,I)\) if and
only if it belongs to
\(\noderef{TwoBitePreTerminalRecordedPairs}(\omega,\varepsilon,I)\), which is
equivalent to membership in \(R\).
\end{proof}
